A healthcare transport network in the Santiago Metropolitan Region must decide every week how many vehicles to hire, where to station them and which transfer hubs to open, and every day how to serve between 455 and 731 weekly requests for clinical samples, refrigerated drugs, supplies and equipment. This report proposes a matheuristic that coordinates both levels. GRASP and VND heuristics generate feasible daily routes, a set-partitioning model solved with Gurobi selects the fleet, bases, hubs and service mode of each request, and a dual-based column generation adds routes that can lower the cost. Compared with a rule-based baseline, the matheuristic is cheaper in all 80 available weeks, saving 9.9\,\% on the large instances (3.9 million CLP per week) and 15.7\,\% on the small ones, mostly by routing 16\,\% fewer kilometres with almost the same fleet. The savings are statistically significant, stable across random seeds, and hold in every sensitivity scenario. Refrigerated vans are the binding resource, and splitting requests across vehicles is essential, while transfer hubs add little value. No base is indispensable and an additional base does not pay off, whereas two extra vans in high-demand weeks reduce cost and outsourcing. A staged implementation with a shadow run and a zone pilot is proposed.

\vspace{0.5em}
\noindent\textbf{Keywords:} vehicle routing, fleet sizing, matheuristic, column generation, healthcare logistics.
